% LỜI TỰA --- Quyển VII
\chapter*{Lời Tựa}
\addcontentsline{toc}{chapter}{Lời Tựa}
\markboth{Lời Tựa}{Lời Tựa}

\begin{truyen}{Lời Tựa}{Teachings of the Wise}
\chuhoa{G}{iáo viên bắt đầu bộ sách này với một câu hỏi đơn giản: làm thế nào để mỗi buổi học có thể để lại thứ gì đó đáng nhớ hơn là bài thi?}
\textit{(A teacher began this series with a simple question: how can each class leave behind something more memorable than a test?)}

Những bậc hiền nhân không phải là những người không bao giờ sai~--- họ là những người học được từ sai lầm của mình và của người khác.
\textit{(The wise are not those who never make mistakes~--- they are those who learn from their own mistakes and others'.)}

Quyển bảy tập hợp lời dạy từ nhiều bậc thầy qua các thời đại: triết gia Đông Tây, nhà giáo dục, người kinh doanh, người sống đến tuổi già với trí tuệ tinh tường. Những lời này không có tuổi~--- một câu từ Khổng Tử hay Seneca vẫn còn nguyên giá trị trong lớp học hôm nay.
\textit{(Volume seven gathers teachings from great teachers across the ages: Eastern and Western philosophers, educators, businesspeople, and those who lived to old age with clear wisdom. These words are timeless~--- a saying from Confucius or Seneca still holds full value in today's classroom.)}

Bộ sách <<Truyện Tích Cảm Hứng Song Ngữ>> gồm 28 quyển, mỗi quyển một chủ đề riêng. Quyển VII này mang chủ đề: \textbf{triết lý từ bậc thầy qua các thời đại Đông~-- Tây}.
\textit{(The series ``Inspiring Bilingual Tales'' contains 28 volumes, each with its own theme. Volume VII focuses on: \textbf{triết lý từ bậc thầy qua các thời đại Đông~-- Tây}.})
\end{truyen}

\begin{baihoc}
Đọc lời của bậc hiền nhân không phải để thuộc lòng~--- mà để sống theo.
\textit{(Reading the words of the wise is not to memorize them~--- it is to live by them.)}
\end{baihoc}

\begin{tcolorbox}[
  colback=truyenblue!6, colframe=truyenblue,
  coltitle=white, fonttitle=\bfseries,
  title={Easy English},
  arc=4pt, boxrule=1pt, breakable, enhanced
]
\textbf{Short English to remember:}

Wisdom from the past lights the path ahead.

The greatest teachers speak across centuries.
\end{tcolorbox}

\begin{tcolorbox}[
  colback=truyengold!10, colframe=truyengold!70!black,
  coltitle=black, fonttitle=\bfseries,
  title={T\textit{ự} H\textit{ọ}c Nhanh},
  arc=4pt, boxrule=1pt, breakable, enhanced
]
1. Lời dạy nào trong quyển này em muốn viết lên bảng lớp? \textit{(Which teaching in this volume would you write on the classroom board?)}

2. Những bậc hiền nhân Đông và Tây có điểm gì giống nhau trong lời dạy của họ? \textit{(What do Eastern and Western wise teachers have in common in their teachings?)}

3. Hôm nay em có thể áp dụng ngay một lời dạy nào từ quyển này không? \textit{(Is there a teaching in this volume you can apply today?)}
\end{tcolorbox}

\begin{center}
\textit{\color{truyengold}``Lời dạy của người xưa vẫn soi sáng những bước đi hôm nay.''}
\end{center}

\begin{flushright}
\textbf{Tôi Nguyễn Văn Sang}\\
\textit{GV THPT Nguyễn Hữu Cảnh - TP HCM}
\end{flushright}

\ngancach
